\section{第11章\quad 曲线积分}
\subsection{对弧长的曲线积分}
\begin{exercise}{392}{基础·常数弧长积分}
设曲线 $L$ 长度为 $a>0$。求 $\int_L3\,\dd s$，说明曲线反向是否改变结果。
\end{exercise}
\begin{exercise}{393}{基础·弧长积分非负性}
设 $f$ 在分段光滑曲线 $L$ 上连续且非负。由定义说明 $\int_Lf\,\dd s\geq0$。
\end{exercise}
\begin{exercise}{394}{基础·积分比较}
设曲线 $L$ 长度为 $2$，连续函数满足 $1\leq f\leq3$。给出 $\int_L f\,\dd s$ 的上下界。
\end{exercise}
\begin{exercise}{395}{基础·积分绝对值不等式}
设 $f$ 在分段光滑曲线 $L$ 上连续。利用 $-|f|\leq f\leq|f|$ 证明 $|\int_Lf\,\dd s|\leq\int_L|f|\,\dd s$。
\end{exercise}
\begin{exercise}{396}{基础·圆弧径向函数}
设 $L$ 为半径 $2$ 的下半圆 $x^2+y^2=4,y\leq0$。计算 $\int_L\sqrt{x^2+y^2}\,\dd s$。
\end{exercise}
\begin{exercise}{397}{基础·线段参数}
设 $L$ 连接 $(-1,0)$ 与 $(0,2)$。计算 $\int_L(2x-y+5)\,\dd s$。
\end{exercise}
\begin{exercise}{398}{基础·空间圆周}
设 $L:x^2+y^2=4,z=2$。计算 $\oint_L\arctan\frac{z^2}{x^2+y^2}\,\dd s$。
\end{exercise}
\begin{exercise}{399}{常规·大圆对称性}
设 $L:x^2+y^2+z^2=9,x+y+z=0$。利用三个坐标的对称性及球面方程求 $\oint_Lz^2\,\dd s$。
\end{exercise}
\begin{exercise}{400}{基础·椭圆约束}
设椭圆 $L:x^2/4+y^2/3=1$ 的周长为 $a$。利用对称性及曲线方程计算 $\oint_L(2xy+3x^2+4y^2)\,\dd s$，用 $a$ 表示结果。
\end{exercise}
\begin{exercise}{401}{基础·折线积分}
设 $L:y=|x|,-1\leq x\leq1$。计算 $\int_L(x+y)\,\dd s$。
\end{exercise}
\begin{exercise}{402}{基础·线段弧长平均}
设 $L$ 连接 $(1,0)$ 与 $(0,1)$。求 $\int_Lx^2\,\dd s$ 及 $x^2$ 的弧长平均值。
\end{exercise}
\begin{exercise}{403}{常规·互换对称性}
设 $L:x^2+y^2=1$。利用 $x,y$ 互换计算 $\oint_L\frac{e^x}{e^x+e^y}\,\dd s$。
\end{exercise}
\begin{exercise}{404}{基础·圆周奇函数消去}
设 $L:x^2+y^2=1,z=0$。计算 $\oint_L(xy+x^2+z)\,\dd s$。
\end{exercise}
\begin{exercise}{405}{基础·指数曲线弧长因子}
设 $L:y=e^x,0\leq x\leq1$。计算 $\int_L\sqrt{1+y^2}\,\dd s$。
\end{exercise}
\begin{exercise}{406}{基础·抛物线弧长积分}
设 $L:y^2=2x,0\leq y\leq\sqrt3$。以 $y$ 为参数计算 $\int_Ly\,\dd s$。
\end{exercise}
\begin{exercise}{407}{常规·双曲线弧积分}
设 $L$ 为 $xy=1$ 上 $1/2\leq x\leq1$ 的弧。计算 $\int_L\frac{x+y}{\sqrt{x^2+y^2}}\,\dd s$。
\end{exercise}
\begin{exercise}{408}{基础·平移圆弧参数}
设 $L:x=1+\cos t,y=\sin t,0\leq t\leq\pi/2$。求 $\int_Lxy\,\dd s$。
\end{exercise}
\begin{exercise}{409}{基础·圆心与平均值}
设 $L:(x-1)^2+(y-1)^2=1$。利用对称性求 $\oint_L(x+y)\,\dd s$。
\end{exercise}
\begin{exercise}{410}{基础·空间螺旋弧}
设 $L:x=\cos t,y=\sin t,z=t,0\leq t\leq2\pi$。计算 $\int_Lz\,\dd s$。
\end{exercise}
\begin{exercise}{411}{基础·对数曲线}
设 $L:y=\ln x,1\leq x\leq\sqrt3$。求 $\int_L\frac{e^y}{1+x^2}\,\dd s$。
\end{exercise}
\begin{exercise}{412}{常规·椭圆参数与根式}
设 $L:x^2+4y^2=4,y\geq0$。用 $x=2\cos t,y=\sin t$ 计算 $\int_L\sqrt{1+3y^2}\,\dd s$。
\end{exercise}
\begin{exercise}{413}{基础·极坐标曲线}
设 $L$ 的极坐标方程为 $r=2$，$0\leq\theta\leq\pi/2$。求 $\int_L(x-y)\,\dd s$。
\end{exercise}
\begin{exercise}{414}{基础·空间线段积分}
设 $L$ 连接 $(0,0,0)$ 与 $(1,2,2)$。计算 $\int_L(x+y+z)\,\dd s$。
\end{exercise}
\begin{exercise}{415}{基础·水平截圆}
设 $L:x^2+y^2+z^2=4,z=1$。求圆的半径并计算 $\oint_L(x^2+y^2)\,\dd s$。
\end{exercise}
\begin{exercise}{416}{常规·曲边完整边界}
设 $L$ 为 $y=x^2$ 与 $y=x$ 围成区域的完整边界。分两段计算 $\oint_Lx\,\dd s$。
\end{exercise}
\begin{exercise}{417}{基础·三角形边界}
设 $L$ 为顶点 $(0,0),(1,0),(0,1)$ 的三角形边界。分边计算 $\oint_L(x+y)\,\dd s$。
\end{exercise}
\begin{exercise}{418}{基础·圆周弧长积分变式}
设 $L:x^2+y^2=4$。计算 $\oint_L(x^2+y^2)\,\dd s$。
\end{exercise}
\begin{exercise}{419}{常规·圆形细丝质量与质心}
设 $L:x^2+y^2=1$，线密度为 $\rho=1+x/2$。求质量与质心坐标。
\end{exercise}
\subsection{对坐标的曲线积分}
\begin{exercise}{420}{基础·竖直线段}
设 $L$ 从 $(1,0)$ 到 $(1,2)$。分别计算 $\int_Ly\,\dd x$ 与 $\int_Ly\,\dd y$。
\end{exercise}
\begin{exercise}{421}{基础·常向量场}
设 $L$ 为从 $(1,2)$ 到 $(3,-1)$ 的分段光滑曲线。计算 $\int_L2\,\dd x-3\,\dd y$，并求反向积分。
\end{exercise}
\begin{exercise}{422}{基础·根式与绝对值}
设 $L:y=\sin x$，$x$ 从 $0$ 变到 $\pi$。计算 $\int_L\sqrt{1-y^2}\,\dd x$，提示根式等于 $|\cos x|$。
\end{exercise}
\begin{exercise}{423}{基础·抛物线坐标积分}
设 $L:y=x^2$，从 $(-1,1)$ 到 $(1,1)$。计算 $\int_Lx\,\dd y$。
\end{exercise}
\begin{exercise}{424}{基础·线段与方向}
设 $L$ 为从 $(3,0)$ 到 $(0,1)$ 的线段。计算 $\int_Lxy\,\dd y$。
\end{exercise}
\begin{exercise}{425}{基础·双曲线参数}
设 $L:xy=1$，从 $(1,1)$ 到 $(2,1/2)$。计算 $\int_Lx\ln y\,\dd y$。
\end{exercise}
\begin{exercise}{426}{基础·曲线约束代换}
设 $L:y=\cos x$，从 $(0,1)$ 到 $(\pi,-1)$。利用 $\dd y=-\sin x\,\dd x$ 计算 $\int_Le^y\sin x\,\dd x+(1-y)\,\dd y$。
\end{exercise}
\begin{exercise}{427}{基础·以纵坐标作参数}
设 $L:x=y^2$，从 $(1,-1)$ 到 $(1,1)$。计算 $\int_Ly\,\dd x+x\,\dd y$。
\end{exercise}
\begin{exercise}{428}{常规·摆线参数}
设 $L:x=t-\sin t,y=1-\cos t$，$t:0\to2\pi$。计算 $\int_Ly\,\dd x$。
\end{exercise}
\begin{exercise}{429}{基础·螺旋线坐标积分}
设 $L:x=2\cos t,y=2\sin t,z=t$，$t:0\to\pi$。计算 $\int_Ly\,\dd x+z\,\dd z$。
\end{exercise}
\begin{exercise}{430}{基础·全微分端点法}
设 $L$ 从 $(2,0)$ 沿椭圆 $x^2+4y^2=4$ 的第一象限弧到 $(0,1)$。利用全微分计算 $\int_Ly\,\dd x+x\,\dd y$。
\end{exercise}
\begin{exercise}{431}{基础·角形式圆周积分}
设 $L:x^2+y^2=1$ 取顺时针方向。用参数计算 $\oint_L\frac{y\,\dd x-x\,\dd y}{x^2+y^2}$。
\end{exercise}
\begin{exercise}{432}{基础·圆周约束化简}
设 $L:(x-1)^2+(y-2)^2=9$ 取逆时针方向。利用曲线方程计算 $\oint_L\frac{\dd x+\dd y}{x^2+y^2-2x-4y}$。
\end{exercise}
\begin{exercise}{433}{基础·空间全微分}
设 $L$ 从 $(1,2,3)$ 到 $(-1,4,0)$。求 $\int_Lx\,\dd x+y\,\dd y+z\,\dd z$，写出原函数。
\end{exercise}
\begin{exercise}{434}{基础·给定空间参数}
设 $L:x=\cos t,y=\sin t,z=1-\cos t$，$t:0\to2\pi$。计算 $\oint_Ly\,\dd x+z\,\dd y$。
\end{exercise}
\begin{exercise}{435}{基础·折线分段积分}
设 $L:y=1-|x|$，从 $(-1,0)$ 到 $(1,0)$。分段计算 $\int_Ly\,\dd x+x\,\dd y$。
\end{exercise}
\begin{exercise}{436}{基础·三角形边界积分}
设 $L$ 依次经过 $(0,0),(1,0),(1,1),(0,0)$。分边求 $\oint_Ly\,\dd x+x^2\,\dd y$。
\end{exercise}
\begin{exercise}{437}{基础·复合路径与原函数}
设 $L$ 从 $(0,0)$ 沿任意分段光滑路径到 $(0,2)$。找出原函数并计算 $\int_Ly\,\dd x+(x-1)\,\dd y$。
\end{exercise}
\begin{exercise}{438}{基础·曲边闭合路径}
设 $L$ 从原点沿 $y=\sqrt x$ 到 $(1,1)$，再沿 $y=x$ 回到原点。计算 $\oint_Ly\,\dd x$。
\end{exercise}
\begin{exercise}{439}{基础·空间折线积分}
设 $L$ 依次经过 $(0,0,0),(0,0,2),(1,0,2),(1,3,2)$。计算 $\int_Lz^2\,\dd x+x^2\,\dd y+y^2\,\dd z$。
\end{exercise}
\begin{exercise}{440}{常规·空间曲线平面化}
设 $L:x^2+y^2=1,z=2-y$，从正 $z$ 轴向下看为逆时针。用参数计算 $\oint_Lz\,\dd x+x\,\dd y$。
\end{exercise}
\begin{exercise}{441}{基础·乘积全微分}
设 $L$ 从 $(1,1,1)$ 到 $(2,1,3)$。利用 $\dd(xyz)$ 求 $\int_Lyz\,\dd x+zx\,\dd y+xy\,\dd z$。
\end{exercise}
\begin{exercise}{442}{基础·切向量与两类积分}
设 $L:x=\cos t,y=\sin t$，$t:0\to\pi/2$。写出同向单位切向量，将 $\int_LP\,\dd x+Q\,\dd y$ 化为对弧长的积分，其中 $P,Q$ 连续。
\end{exercise}
\begin{exercise}{443}{基础·参数反向}
设 $L:x=t,y=t^2,z=t^3$，$t:1\to0$。计算 $\int_L\dd x+2y\,\dd y+3z^2\,\dd z$。
\end{exercise}
\begin{exercise}{444}{常规·抛物线积分转换}
设 $L:y=x^2$，从 $(1,1)$ 到 $(0,0)$。计算 $\int_L\frac{\dd x+2x\,\dd y}{1+4x^2}$，并写出该积分对应的参数积分。
\end{exercise}
\begin{exercise}{445}{基础·原路返回变式}
设 $C$ 从 $(0,0)$ 沿线段到 $(1,0)$，再原路返回。分别计算 $\int_C1\,\dd s$ 与 $\int_C1\,\dd x$。
\end{exercise}
\begin{exercise}{446}{常规·两条路径小综合}
分别沿 $y=x$ 与 $y=x^2$ 从 $(0,0)$ 到 $(1,1)$ 计算 $\int_Ly\,\dd x$，据此判断该积分是否与路径无关。
\end{exercise}
\subsection{格林公式及其应用}
\subsubsection{封闭曲线、补线与奇点}
\begin{exercise}{447}{基础·闭曲线全微分}
设 $L$ 为分段光滑闭曲线。找出原函数并计算 $\oint_L(y-\cos x)\,\dd x+(x+\sin y)\,\dd y$。
\end{exercise}
\begin{exercise}{448}{基础·环流与面积}
设 $D$ 的边界分段光滑且取正向，已知 $\oint_{\partial D}(1+3y)\,\dd x+(5x-1)\,\dd y=4$。用格林公式求 $D$ 的面积。
\end{exercise}
\begin{exercise}{449}{基础·负向边界}
设 $L$ 为顶点 $(0,0),(3,0),(3,2)$ 的三角形顺时针边界。用格林公式求 $\oint_Ly\,\dd x+5x\,\dd y$。
\end{exercise}
\begin{exercise}{450}{基础·曲边区域格林公式}
设 $L$ 为 $D=\{0\leq x\leq1,x^2\leq y\leq x\}$ 的顺时针边界。计算 $\oint_Ly\,\dd x-x\,\dd y$。
\end{exercise}
\begin{exercise}{451}{基础·圆周环流}
设 $L:x^2+y^2=1$ 取逆时针方向。用格林公式求 $\oint_Lx^2y\,\dd x-xy^2\,\dd y$。
\end{exercise}
\begin{exercise}{452}{基础·扇环边界}
设 $D=\{(r\cos\theta,r\sin\theta):1\leq r\leq2,0\leq\theta\leq\pi/2\}$。计算完整正向边界上的 $\oint_{\partial D}-y\,\dd x+x\,\dd y$。
\end{exercise}
\begin{exercise}{453}{基础·抛物线弓形}
设 $L$ 为 $y=x^2$ 与 $y=1$ 围成区域的逆时针边界。用格林公式求 $\oint_L(e^x+y)\,\dd x-x\,\dd y$。
\end{exercise}
\begin{exercise}{454}{基础·菱形边界}
设 $L$ 为菱形 $|x|+|y|\leq1$ 的逆时针边界。计算 $\oint_L-y\,\dd x+2x\,\dd y$。
\end{exercise}
\begin{exercise}{455}{基础·折线补线}
设 $L$ 依次连接 $(-1,0),(0,1),(1,0)$。补上 $x$ 轴线段，用格林公式计算 $\int_Ly\,\dd x$。
\end{exercise}
\begin{exercise}{456}{常规·半圆弧补线}
设 $L:x^2+y^2=1,y\geq0$，从 $(-1,0)$ 到 $(1,0)$。用补线法计算 $\int_L(x^2-3y)\,\dd x+x\,\dd y$。
\end{exercise}
\begin{exercise}{457}{基础·曲线方程化简}
设 $L:x=y^2$，从 $(1,-1)$ 到 $(1,1)$。计算 $\int_L\frac{y\,\dd x+x\,\dd y}{y^2-x+2}$。
\end{exercise}
\begin{exercise}{458}{常规·带孔区域格林公式}
设 $D=\{1\leq x^2+y^2\leq4\}$。计算完整正向边界上的 $\oint_{\partial D}\frac{-y\,\dd x+x\,\dd y}{x^2+y^2}$，写明内外圆周的方向及格林公式的适用条件。
\end{exercise}
\begin{exercise}{459}{基础·避开奇点的圆周}
设 $L:(x-2)^2+y^2=1$ 取逆时针方向。用格林公式计算 $\oint_L\frac{-y\,\dd x+x\,\dd y}{x^2+y^2}$，说明原点是否在闭圆盘内。
\end{exercise}
\subsubsection{路径无关、原函数与全微分方程}
\begin{exercise}{460}{基础·保守场参数}
求使 $(ax+y)\,\dd x+(x+2y)\,\dd y$ 在平面内为全微分的参数 $a$，并写出一个原函数。
\end{exercise}
\begin{exercise}{461}{基础·系数匹配}
求使 $(2xy+a)\,\dd x+(bx^2+2y)\,\dd y$ 在平面内与路径无关的 $b$，其中 $a$ 为常数；求一个原函数。
\end{exercise}
\begin{exercise}{462}{基础·路径无关条件}
检验 $e^{xy}(y\,\dd x+x\,\dd y)$ 是否为全微分，并求从 $(0,0)$ 到 $(1,1)$ 的积分。
\end{exercise}
\begin{exercise}{463}{基础·分式原函数}
在半平面 $x+y>0$ 内，求 $\frac{y\,\dd x-x\,\dd y}{(x+y)^2}$ 的一个原函数。提示：考虑 $x/(x+y)$。
\end{exercise}
\begin{exercise}{464}{基础·复合函数全微分}
计算从 $(1,0)$ 到 $(2,1)$ 的 $\int_L\frac{y\,\dd x+x\,\dd y}{1+x^2y^2}$，写出原函数。
\end{exercise}
\begin{exercise}{465}{常规·反求场分量}
设 $f\in C^1(\mathbb R)$，$f(0)=1$，且 $2xy\,\dd x+f(x)\,\dd y$ 在平面内与路径无关。求 $f$ 及一个原函数。
\end{exercise}
\begin{exercise}{466}{基础·定义域条件}
在半平面 $x>0$ 内求 $\frac{-y\,\dd x+x\,\dd y}{x^2+y^2}$ 的原函数，并求从 $(1,0)$ 到 $(1,1)$ 的积分。
\end{exercise}
\begin{exercise}{467}{基础·乘积原函数}
找出 $(2xy+\cos x)\,\dd x+x^2\,\dd y$ 的原函数，求从 $(0,0)$ 到 $(1,2)$ 的积分。
\end{exercise}
\begin{exercise}{468}{基础·端点消去}
设 $L$ 从 $(0,0)$ 到 $(2,0)$。计算 $\int_L(e^x+2xy)\,\dd x+(x^2+e^{-y^2})\,\dd y$。原函数中允许保留一元定积分。
\end{exercise}
\begin{exercise}{469}{基础·多项式原函数}
计算从 $(1,2)$ 到 $(-1,1)$ 的 $\int_L(2x-y)\,\dd x+(2y-x)\,\dd y$，写出原函数。
\end{exercise}
\begin{exercise}{470}{基础·径向指数场}
设 $L$ 从 $(0,0)$ 到 $(1,1)$。利用全微分计算 $\int_Le^{x^2+y^2}(x\,\dd x+y\,\dd y)$。
\end{exercise}
\begin{exercise}{471}{基础·线性场原函数}
证明 $(x+2y)\,\dd x+(2x-y)\,\dd y$ 在平面内与路径无关，计算从 $(-1,-2)$ 到 $(3,4)$ 的积分。
\end{exercise}
\begin{exercise}{472}{基础·混合项配对}
求 $(xy^2+x)\,\dd x+(x^2y+y)\,\dd y$ 的原函数，并计算从 $(0,0)$ 到 $(1,1)$ 的积分。
\end{exercise}
\begin{exercise}{473}{基础·椭圆面积}
设 $a,b>0$，$L:x=a\cos t,y=b\sin t,t:0\to2\pi$。用 $\frac12\oint_L(x\,\dd y-y\,\dd x)$ 求椭圆面积。
\end{exercise}
\begin{exercise}{474}{常规·星形线面积}
设 $L:x=\cos^3t,y=\sin^3t,t:0\to2\pi$。用曲线积分求所围区域面积。
\end{exercise}
\begin{exercise}{475}{基础·归一化原函数}
求满足 $u(0,0)=0$ 且 $\dd u=(2x+\sin y)\,\dd x+x\cos y\,\dd y$ 的 $u(x,y)$。
\end{exercise}
\begin{exercise}{476}{基础·半平面原函数}
求 $\frac{y\,\dd x-x\,\dd y}{x^2+y^2}$ 在 $x>0$ 内满足 $u(1,0)=0$ 的原函数。
\end{exercise}
\begin{exercise}{477}{常规·参数与初值}
设 $\dd u=(x+a)e^{x+y}\,\dd x+xe^{x+y}\,\dd y$。求 $a$ 以及满足 $u(0,0)=1$ 的原函数。
\end{exercise}
\begin{exercise}{478}{基础·全微分方程}
检验 $(2x+y)\,\dd x+(x+2y)\,\dd y=0$ 为全微分方程，并求隐式通解。
\end{exercise}
\begin{exercise}{479}{基础·全微分方程初值}
求全微分方程 $(e^x+y)\,\dd x+(x+2y)\,\dd y=0$ 经过 $(0,1)$ 的积分曲线方程。
\end{exercise}
\begin{exercise}{480}{基础·原函数计算变式}
求 $2xy\,\dd x+(x^2+3y^2)\,\dd y$ 的一个原函数，并用偏导数检验。
\end{exercise}
\begin{exercise}{481}{常规·路径无关小综合}
设 $\omega=(ay+x)\,\dd x+(2x+y)\,\dd y$。求使它在整个平面内与路径无关的 $a$；在此时求从 $(0,0)$ 到 $(1,2)$ 的积分。
\end{exercise}
\subsection{曲线积分的应用}
\begin{exercise}{482}{基础·变限积分曲线弧长}
设 $L:y=\int_0^x\tan t\,\dd t,0\leq x\leq\pi/4$。求弧长。
\end{exercise}
\begin{exercise}{483}{基础·螺旋线弧长}
设 $L:x=2\cos t,y=2\sin t,z=3t,0\leq t\leq2\pi$。求弧长。
\end{exercise}
\begin{exercise}{484}{基础·抛物线细丝质量}
设 $L:y=x^2,-1\leq x\leq1$，线密度为 $\rho=|x|$。求质量及质心横坐标。
\end{exercise}
\begin{exercise}{485}{基础·空间细丝质量}
设 $L:x=t,y=t^2/2,z=t^3/3,0\leq t\leq1$，线密度为 $\rho=\sqrt{1+2y+3xz}$。求质量。
\end{exercise}
\begin{exercise}{486}{基础·细棒质心}
细棒位于 $x$ 轴的 $[0,1]$ 上，线密度为 $\rho=1+x$。求质量与质心。
\end{exercise}
\begin{exercise}{487}{基础·螺旋细丝重心}
设 $L:x=\cos t,y=\sin t,z=t,0\leq t\leq2\pi$，线密度为 $1$。求质量与重心高度。
\end{exercise}
\begin{exercise}{488}{基础·半圆细丝惯量}
设 $L:x^2+y^2=1,y\geq0$，线密度为 $1$。求关于 $x$ 轴和 $y$ 轴的转动惯量。
\end{exercise}
\begin{exercise}{489}{基础·圆形细丝轴惯量}
设 $L:x^2+y^2=4,z=0$，线密度为 $1$。求质量与关于 $z$ 轴的转动惯量。
\end{exercise}
\begin{exercise}{490}{基础·变力做功}
质点沿 $x=\cos t,y=\sin t,z=t$，$t:0\to\pi$，在力 $\bm F=(y,-x,z)$ 作用下移动。求力做的功。
\end{exercise}
\begin{exercise}{491}{基础·回复力做功}
力为 $\bm F=(-2x,-2y)$。求质点从 $(2,0)$ 沿任意分段光滑路径到 $(0,1)$ 时力做的功。
\end{exercise}
\begin{exercise}{492}{常规·径向力做功}
设 $\bm F=(x,y,z)/(x^2+y^2+z^2)$，质点沿线段从 $(1,0,0)$ 到 $(2,0,0)$。求功，写明路径上分母不为零。
\end{exercise}
\begin{exercise}{493}{基础·细丝平行轴公式}
设细丝质量为 $M>0$，质心纵坐标为 $\bar y$，线密度非负。展开 $(y-c)^2$，证明关于直线 $y=c$ 的惯量满足 $I(c)=I(\bar y)+M(c-\bar y)^2$。
\end{exercise}
\begin{exercise}{494}{常规·半圆质心与做功}
单位上半圆细丝的线密度为 $1$。求质心；再求力 $\bm F=(-y,x)$ 沿上半圆从 $(1,0)$ 逆时针到 $(-1,0)$ 所做的功。
\end{exercise}
\subsubsection{常规综合练习}
\begin{exercise}{495}{基础·边界互换对称性}
设 $L$ 为 $y=x^2$ 与 $x=y^2$ 在第一象限围成区域的完整边界。利用 $x,y$ 互换求 $\oint_L(x-y)\,\dd s$。
\end{exercise}
\begin{exercise}{496}{基础·圆周象限分解}
设 $L:x^2+y^2=1$。计算 $\oint_L(|x|+|y|)\,\dd s$。
\end{exercise}
\begin{exercise}{497}{基础·悬链线弧长因子}
设 $L:y=(e^x+e^{-x})/2,-1\leq x\leq1$。求 $\int_L|x|\,\dd s$。提示：$\sqrt{1+(y')^2}=(e^x+e^{-x})/2$。
\end{exercise}
\begin{exercise}{498}{基础·空间三角形弧长积分}
设 $L$ 为顶点 $(0,0,0),(1,0,0),(1,0,1)$ 的三角形完整边界。分边计算 $\oint_Lz\,\dd s$。
\end{exercise}
\begin{exercise}{499}{基础·扇形完整边界}
设 $L$ 为 $0\leq r\leq1,0\leq\theta\leq\pi/2$ 的扇形完整边界。分为两条半径和圆弧，计算 $\oint_L\sqrt{x^2+y^2}\,\dd s$。
\end{exercise}
\begin{exercise}{500}{基础·等值曲线积分}
设 $L$ 为 $f(x,y)=1$ 的等值线上从 $(0,1)$ 到 $(1,0)$ 的光滑弧，长度为 $\ell$，$f\in C^1$。计算 $\int_Lf\,\dd x$ 与 $\int_Lf\,\dd s$。
\end{exercise}
\begin{exercise}{501}{常规·曲边区域环流}
设 $L$ 为 $y=x^2$ 与 $y=\sqrt x$ 围成区域的顺时针边界。用格林公式计算 $\oint_L-y\,\dd x+x\,\dd y$。
\end{exercise}
\begin{exercise}{502}{基础·正弦曲线做功}
设 $L:y=\sin x$，从 $(0,0)$ 到 $(\pi,0)$。计算 $\int_Ly\,\dd x+2x\,\dd y$。
\end{exercise}
\begin{exercise}{503}{基础·空间闭合曲线}
设 $L:x=\cos t,y=\sin t,z=1$，$t:0\to2\pi$。计算 $\oint_Ly^2\,\dd x+z^2\,\dd y+x^2\,\dd z$。
\end{exercise}
\begin{exercise}{504}{基础·积分估值}
设有向曲线 $L$ 长度为 $3$，连续函数满足 $P^2+Q^2\leq4$。利用曲线积分估值不等式给出 $|\int_LP\,\dd x+Q\,\dd y|$ 的上界。
\end{exercise}
\begin{exercise}{505}{常规·线性变换与面积}
设 $L:|x|+|x+y|=1$ 取逆时针方向。用 $u=x,v=x+y$ 求包围面积，再计算 $\oint_L(x-y)\,\dd x+(x+y)\,\dd y$。
\end{exercise}
\begin{exercise}{506}{基础·含函数场的格林公式}
设 $z(x,y)=x+y$，$L$ 为逆时针单位圆周。计算 $\oint_L(z-3y)\,\dd x+(z+x^2)\,\dd y$。
\end{exercise}
\begin{exercise}{507}{基础·线性场环流}
设 $L:(x-1)^2+(y-1)^2=1$ 取逆时针方向。用格林公式求 $\oint_L-2y\,\dd x+3x\,\dd y$。
\end{exercise}
\begin{exercise}{508}{常规·余弦弧补线}
设 $L:y=\cos x$，从 $(-\pi/2,0)$ 到 $(\pi/2,0)$。用补线法计算 $\int_L(e^x-y)\,\dd x+x\,\dd y$。
\end{exercise}
\begin{exercise}{509}{基础·圆盘环流}
设 $L:x^2+y^2=4$ 取逆时针方向。用格林公式求 $\oint_L(1-x^2/4)y\,\dd x$。
\end{exercise}
\begin{exercise}{510}{基础·路径无关条件检验}
设 $f\in C^1(\mathbb R)$，$f(0)=0$，$f(x)\,\dd x+(x+y)\,\dd y$ 是否可能在平面内与路径无关？写出交叉偏导数并说明结论。
\end{exercise}
\begin{exercise}{511}{基础·同端点原函数法}
设 $L$ 为从 $(0,0)$ 到 $(1,1)$ 的任意分段光滑曲线。求 $\int_L(1-2xy)\,\dd x-x^2\,\dd y$。
\end{exercise}
\begin{exercise}{512}{常规·反求未知分量}
设 $f\in C^1(0,\infty)$，$f(1)=0$，$yf(x)\,\dd x+(x^2-2x)\,\dd y$ 在 $x>0$ 内与路径无关。求 $f$ 及一个原函数。
\end{exercise}
\begin{exercise}{513}{基础·由原函数求场分量}
已知原函数为 $u=x^2y+y^2$。求 $P,Q$ 使 $\dd u=P\,\dd x+Q\,\dd y$，并求从 $(0,0)$ 到 $(1,2)$ 的积分。
\end{exercise}
\begin{exercise}{514}{常规·偏导与端点积分}
设 $f_x=2x+y$，$f(0,y)=y^2$。求 $f$，并计算从 $(0,0)$ 到 $(1,2)$ 的 $\int_Lf_x\,\dd x+f_y\,\dd y$。
\end{exercise}
\begin{exercise}{515}{基础·三角形面积积分}
设 $L$ 依次经过 $(0,0),(2,0),(0,3),(0,0)$。计算 $\frac12\oint_L(x\,\dd y-y\,\dd x)$，说明其与三角形面积的关系。
\end{exercise}
\begin{exercise}{516}{常规·全微分方程参数}
求使 $(2xy+a)\,\dd x+(bx^2+2y)\,\dd y=0$ 为全微分方程的 $b$，其中 $a$ 为常数；求经过 $(0,0)$ 的积分曲线方程。
\end{exercise}
\begin{exercise}{517}{基础·弧长表达式}
分别写出 $y=\sin x$（$0\leq x\leq\pi$）的弧长积分，以及椭圆 $x=2\cos t,y=\sin t$（$0\leq t\leq2\pi$）的周长积分，无需求值。
\end{exercise}
\begin{exercise}{518}{基础·细棒质心变式}
设细棒位于 $x$ 轴的 $[0,2]$，线密度为 $\rho=x$。求质量与质心。
\end{exercise}
\begin{exercise}{519}{常规·旋转力场小综合}
设 $\bm F=(-y,x)$。分别求质点从 $(1,0)$ 沿单位圆上半圆逆时针到 $(-1,0)$，以及沿直径线段到同一终点时力所做的功，据此判断路径是否影响功。
\end{exercise}
